\chapter{Discussion}
\label{chap:discussion}

This chapter discusses the design decisions made in this thesis and reflects on their implications for extensible tool support in OCPM.
It also highlights limitations of the current approach and outlines open challenges that remain for future research and development.

The main goal of this thesis was to reduce fragmentation in the current OCPM tool landscape.
The analysis of existing tools showed that many systems reimplement similar functionality, have complex installation procedures, or lack long-term maintainability.
The requirements derived directly from these observations proved highly useful in guiding the design of \textit{Ocelescope}.

Implementing \textit{Ocelescope} as a web application makes it possible to run the interface in a responsive and cross-platform environment.
The tool can run locally in a browser or be hosted.
Running \textit{Ocelescope} inside a container allows the creation of images that remain stable despite changing dependencies.
These images can run on all major operating systems and only require a container engine.
Containers also provide isolated environments that increase protection during execution.

Using Docker to build and manage containers simplifies installation.
Users can run \textit{Ocelescope} with a single command, and predefined startup scripts set up configuration and enable distribution through an image registry.
This registry also allows the creation of versioned images.
The tool evaluation showed that Docker is still manageable for users with some technical experience.
It is less clear how non-technical users would set it up.
Since the module system targets domain users, hosting \textit{Ocelescope} remains a practical alternative.
Packaging \textit{Ocelescope} into operating-system-dependent installers is possible but should only be considered once the need becomes clear.

Splitting the web application into a TypeScript frontend and a Python backend allows the use of React for the user interface while relying on Python for process mining tasks.
Both parts can then be developed within ecosystems where they are best supported.
This separation also aligns with common development practices in modern web-based tools.

Supporting two extension mechanisms, plugins and modules, allows \textit{Ocelescope} to serve different purposes.
Plugins provide an environment for running OCPM techniques inside a unified tool without extra installation steps.
Modules allow developers to use \textit{Ocelescope} as the foundation for specialized OCPM applications aimed at domain experts.
The presence of both systems increases the complexity of the main framework, but it also ensures that improvements benefit both plugin users and module based tools.

Plugins are designed as a simple way to package existing prototypes so that they run inside the \textit{Ocelescope} framework.
This design encourages integration and, supported by good documentation, can motivate the creation of more OCPM plugins than currently exist for ProM.
Defining plugins as functions with clear inputs and outputs follows established approaches used in other systems.
Automatic generation of input interfaces reduces development effort but also imposes restrictions on possible user interfaces.
Domain-specific inputs, such as graphical constraint builders, cannot be expressed naturally within the existing input system.
Although this is a limitation, developers can still implement such interfaces through modules.
This approach remains easier than building a complete tool from scratch.
More specific input types can be added later if stable patterns emerge, for example a general constraint graph builder.

The resource definitions introduced in this thesis provide an easy way to define artifacts that use an automatic exchange format.
The visualization options are limited, but they provide a solid foundation for most use cases.
If new visualization needs become common, additional types can be added to benefit all plugin developers.

Packaging plugins as simple zip archives keeps development accessible.
This choice restricts implementation to Python and Rust through Python bindings, but this is acceptable at the current stage.
Plugins must run in the backend, which makes it infeasible to host \textit{Ocelescope} as a public plugin execution environment.
Although these limitations are significant, alternative designs are currently not practical and may be reconsidered when underlying technologies mature.

Modules fulfill their role as the place where \textit{Ocelescope} becomes a canvas for creating custom OCPM tools.
They require both frontend and backend components to be built together with the rest of the framework.
This allows better cohesion of the interface and improved backend performance because the backend can access all OCEL data in the session.
Workflows support this process and simplify the build pipeline.
They allow developers to create derivatives without deeper knowledge of the overall system.
Modules allow developers to create tools that match the needs of specific user groups and help reduce fragmentation.
Developers no longer need to redesign core functionality such as log management, inspection, or filtering.
At the same time, they can build new components that benefit all \textit{Ocelescope} tools.

In summary, the design of \textit{Ocelescope} reflects a balance between accessibility, extensibility, and long-term maintainability.
Several trade-offs were unavoidable.
The dual extension system increases complexity but supports different types of users.
The simplified plugin interface limits expressiveness but enables fast development.
The use of containers increases stability but may reduce accessibility for non-technical users.
These decisions were made to create a practical foundation for OCPM tool support and to address the requirements derived from the tool analysis.
The user studies and the proof-of-concept plugins indicate that the key plugin requirements are already met in practice.
End-users were able to install \textit{Ocelescope} and execute plugins, and plugin developers were able to implement and integrate a plugin, with generally positive ratings for perceived usefulness and ease of use.
The feedback from both studies also provided clear directions for refinement, in particular more responsive feedback during plugin execution, faster iteration during development, and clearer guidance for input configuration.
At the same time, requirements related to broader base functionality and embedded extensions were not evaluated comprehensively and should be evaluated in future work.
